% Generated draft: requires mathematical and source review.
\documentclass[11pt]{article}
\usepackage[margin=25mm]{geometry}
\usepackage{fontspec}
\setmainfont{Latin Modern Roman}
\usepackage{amsmath,amssymb,mathtools,unicode-math}
\setmathfont{Latin Modern Math}
\usepackage{enumitem,needspace,xurl,hyperref}
\setlength{\parindent}{0pt}
\setlength{\parskip}{5pt}
\setlength{\emergencystretch}{3em}
\newcommand{\N}{\mathbb{N}}
\newcommand{\Z}{\mathbb{Z}}
\newcommand{\Q}{\mathbb{Q}}
\newcommand{\R}{\mathbb{R}}
\newcommand{\C}{\mathbb{C}}
\newcommand{\F}{\mathbb{F}}
\newcommand{\A}{\mathbb{A}}
\newcommand{\PP}{\mathbb{P}}
\newcommand{\E}{\mathbb{E}}
\newcommand{\eps}{\varepsilon}
\newcommand{\dd}{\,\mathrm{d}}
\begin{document}
\section*{UnsolvedMath ID: 30002058}
\textbf{Source problem code:} catalogue entry
\subsection{Definitions and mathematical statement}
Let a finite commutative ring with identity mean a commutative ring \(R\) whose underlying set is finite and which has a multiplicative identity. The source does not explicitly require \(1_R\ne 0_R\). For such a ring, set
\[
\operatorname{SL}_2(R)=\{A\in M_2(R):\det A=1_R\}.
\]
The notation \(\operatorname{PSL}(2,R)\) below is retained from the source. Its precise quotient convention over general commutative rings is not supplied; in particular, quotienting by \(\{I,-I\}\) and quotienting by all scalar matrices in \(\operatorname{SL}_2(R)\) need not give the same group. This convention must be resolved before the first assertion is fully specified.

The supplied record presents the following two conjectural assertions.
\begin{enumerate}
\item For every compact \(3\)-manifold \(M\) and every \(\gamma\in\pi_1(M)\setminus\{1\}\), there exist a finite commutative ring \(R\) with identity and a group homomorphism
\[
\varphi:\pi_1(M)\longrightarrow\operatorname{PSL}(2,R)
\qquad\text{such that}\qquad
\varphi(\gamma)\ne 1.
\]
The ring and homomorphism may depend on both \(M\) and \(\gamma\). No orientability assumption or exclusion of boundary is stated for this assertion. The notation \(\pi_1(M)\) requires a basepoint; the source does not specify connectedness or how a component is selected when \(M\) is disconnected.
\item For every closed oriented \(3\)-manifold \(M\) not homeomorphic to \(S^3\), there exist a one-vertex triangulation \(T\) of \(M\) and a finite commutative ring \(R\) with identity such that the Thurston equations associated with \(T\) have a solution over \(R\).

The source does not provide the Thurston equations, their variables, their admissibility conditions over rings, or the precise meaning of a one-vertex triangulation. Consequently this second assertion records the supplied target but does not specify a complete algebraic system. In particular, no convention concerning invertible shape parameters, degeneracies, or face identifications is imposed here.
\end{enumerate}
The supplied dated assessment describes Conjecture 2 as weaker, but no precise implication between the two assertions is supplied or asserted here. Both assertions require the unresolved source conventions above to be checked.
\subsection{Short English statement}
The first conjecture asks whether every nonidentity element of a compact \(3\)-manifold group can be detected by a homomorphism to \(\operatorname{PSL}(2,R)\) for some finite commutative ring \(R\) with identity. The second asks whether every closed oriented \(3\)-manifold other than \(S^3\) admits a one-vertex triangulation whose Thurston equations have a solution over such a ring. The projective-group convention and the equation system are unresolved in the supplied evidence.
\subsection{Sources}
\begin{itemize}
\item Supplied local UnsolvedMath registry record, permanent ID \(30002058\), source problem code \texttt{catalogue entry}, title \emph{Finite-Ring Solutions of Thurston Equations}, in the Oberwolfach Reports Open Problems collection. No individual proposer is supplied; the proposed year is recorded as 2012.
\item Supplied report citation: \emph{Triangulations (2013)}, Oberwolfach Reports, DOI 10.4171/owr/2012/24, pp.~1405--1486. Source URL and literature-source link labeled ``original OWR report'': \url{https://doi.org/10.4171/owr/2012/24}. This citation is supplied attribution and has not been independently inspected.
\item Literature assessment dated exactly \texttt{2026-08-21}: the supplied curation labels the problem open and says that the extraction names conjectures and a weaker Conjecture 2 but omits their mathematical statements, preventing a precise claim from being verified or classified as solved or disproved. This is supplied dated curation, not an independently verified literature assessment. The current supplied record does contain two conjectural formulations, leaving the applicability of that earlier assessment unresolved.
\item The registry labels the statement \texttt{corrected\_verified} and reports a manual semantic, typographic, and LaTeX review claiming a self-contained, meaning-preserving cleanup. These are supplied review claims; they do not resolve the missing mathematical conventions identified in this draft.
\end{itemize}
\end{document}
